\label{sec:open}


\newcommand{\highlightpara}[1]{\textbf{\textcolor{openmossblue}{#1}}}



While the emergence of World Action Models (WAMs) marks a pivotal shift from token-level prediction to state-level prediction, the path toward general physical intelligence remains fragmented. We argue that addressing these issues requires more than incremental scaling; it necessitates a fundamental rethinking of how world models should be architectured, grounded, and verified. Consequently, we outline several challenges and opportunities that we believe will define the next phase of WAM research.


\highlightpara{Architectural Coupling.} The field has produced a remarkable diversity of structural strategies for coupling world prediction and action generation---cascaded pipelines, joint diffusion backbones, discrete tokenization schemes, and implicit representation alignment---yet no systematic, controlled study has evaluated these paradigms against one another under matched conditions of scale, data, and evaluation protocol. It remains empirically unclear whether explicit visual prediction is strictly necessary for physical grounding, when cascaded and joint architectures differ in ways that matter for downstream control, and what inductive biases each coupling mechanism actually provides. Rigorous ablation studies and theoretical analysis of the design space are urgently needed to move the field beyond architectural fashion toward principled design choices.



One potentially productive direction lies in questioning whether explicit pixel-space prediction is a necessary component of the coupling altogether. Recent evidence suggests that the primary utility of world modeling in certain WAMs may stem from the auxiliary gradients provided during training rather than the explicit generation of future frames during inference; accordingly, some frameworks demonstrate that removing the future prediction head at test time does not necessarily degrade downstream control performance. This observation opens the door to a more computationally efficient paradigm: instead of reconstructing high-dimensional future observations in pixel space, models could predict abstract representations of future states in a jointly learned latent space. Such a \textit{latent-predictive} approach—exemplified by, but not limited to, JEPA could allow the model to sidestep the pixel-prediction bottleneck by focusing on the causal invariants of the environment while ignoring perceptually irrelevant details. By shifting the objective from pixel-level reconstruction to action-conditioned latent state transitions, WAMs may achieve a more compact and physically grounded coupling between world understanding and action execution, though the extent to which these latent-only representations can maintain the same grounding quality as generative approaches remains a vital open question.


\highlightpara{Multimodal Physical State Representation.} Existing WAMs almost universally predict future world states in the RGB visual modality. However, much of the physical information most critical for contact-rich manipulation is largely invisible in pixel space, including tactile distributions, contact forces, acoustic signatures, and material compliance. A world model confined to visual prediction thus has a systematic blind spot precisely at the physical interactions it most needs to model. Extending WAMs to jointly predict and reason over tactile, force, and proprioceptive future states, while developing the architectures and datasets necessary to support such multimodal world modeling, constitutes an important and largely unexplored frontier.


More broadly, this challenge suggests a shift in how the
\textit{world state} itself is defined within WAMs.
If the state relevant to contact-rich manipulation is not a pixel array but a joint distribution over visual, tactile, and force modalities, then multimodal prediction is not a mere extension of the current paradigm. It demands a reconceptualization of what a world model predicts and for what purpose. A valuable capability for future architectures is, therefore, \textit{modality-adaptive prediction}: the capacity to generate physically grounded predictions when rich sensor streams are available, while degrading gracefully to visual-only inference when they are not, rather than treating full multimodality as a fixed architectural requirement.



\highlightpara{Data Utilization and Mixture Design.} Several works have demonstrated that incorporating egocentric human video and other non-robot data sources can substantially improve downstream manipulation performance, but the principles governing optimal data mixture design remain poorly understood. What is the marginal contribution of each data source as a function of its scale and domain gap? Are the benefits of human video pretraining primarily semantic or dynamical? How should training curricula transition from broad internet-scale priors to precise action-labeled robot demonstrations? Developing principled recipes to these questions is essential for scaling WAM training beyond the current reliance on expensive robot-centric teleoperation data.


We argue that the central challenge in data mixture design lies in disentangling the multiple, potentially overlapping roles that non-robot data plays in grounding. Beyond mere semantic augmentation, we can conceptualize the value of human video through a hierarchy of transferrable knowledge: (1) \textit{low-level physical priors}, such as object permanence and gravitational constraints; (2) \textit{mid-level causal dynamics}, which encode the causal relationship between specific interactions and their physical outcomes; and (3) \textit{high-level task logic}, which encode task-relevant temporal dependencies independent of specific embodiments. Shifting the focus from empirical hyperparameter tuning to a more principled information-theoretic perspective on data mixture could allow the field to identify which specific components of world modeling are best learned from internet-scale video versus precision robot demonstrations. A key research frontier is thus the development of embodiment-aware filtering mechanisms—architectures that can selectively distill universal physical laws from diverse sources while isolating and suppressing behaviors that are kinematically incompatible with the target robot.



\highlightpara{Long-Horizon Planning and Temporal Abstraction.}
WAMs are most commonly evaluated on short-horizon manipulation tasks within single interaction contexts, but genuine embodied generalism demands sustained reasoning over extended task horizons. Current architectures face compounding challenges in this regime: distributional drift in world model predictions accumulates over long rollouts, action errors compound in the absence of corrective replanning, and representing full long-horizon task trajectories as continuous generative outputs is both computationally prohibitive and architecturally unsupported. A principled framework for hierarchical world-action modeling---connecting high-level semantic task decomposition to low-level physical prediction within a unified, learnable architecture---remains a critical open challenge.


We envision three complementary pathways toward long-horizon reasoning. First is the \textit{modular hierarchy}, where WAMs function as low-level physical executors guided by high-level planners (e.g., Vision-Language Models) that decompose complex missions into semantic subgoals. Second is the \textit{intrinsic hierarchical WAM}, an architecture capable of generating multi-resolution future predictions—predicting coarse-grained state transitions for strategic planning while simultaneously synthesizing fine-grained physical details for reactive control. Third is the \textit{scaling of temporal context}, which involves architectural innovations to expand the "memory" of WAMs, allowing them to maintain rich historical state information without the quadratic overhead of standard attention. Whether these paths converge into a unified, end-to-end learnable hierarchy or remain as a composite system of specialized modules is a foundational question for the next generation of WAM research.


\highlightpara{Inference Latency and Computational Efficiency.}
The integration of world prediction imposes a severe ``latency tax'' that threatens the viability of closed-loop control. While early cascaded pipelines suffered from the cumulative delay of independent modules, modern joint WAMs still struggle to meet the millisecond-level response frequencies required for high-fidelity manipulation. DreamZero pushed WAM inference toward the 7Hz through algorithmic acceleration and low-level CUDA optimizations; however, this still lags significantly behind the 50Hz standards of contemporary non-generative VLA policies. The fundamental tension remains:  Can the rich physical foresight of a WAM be preserved at the temporal resolution of real-time motor control?


Underlying this challenge is a theoretical question that has received little direct study: how much predictive fidelity does downstream control actually need? If performance gains level off well before reaching the quality of full diffusion-based synthesis, then the goal should not be to make high-fidelity prediction faster. Instead, it should be to identify the minimal sufficient world model for a given task and act accordingly. This points toward task-adaptive predictive fidelity, where the model dynamically adjusts the depth and resolution of its predictions based on the task's demands and error tolerance， investing computation where precision matters most, and using coarser approximations elsewhere.


\highlightpara{Evaluation Methodology.} Current WAM evaluation suffers from a severe decoupling. World modeling is typically assessed via pixel-level metrics (e.g., PSNR, FVD) that capture visual plausibility but ignore physical correctness---allowing videos where unsupported objects hover or fluids defy gravity to score highly. Conversely, action generation is measured solely by downstream task success. This disjointed paradigm ignores the core premise of WAMs. A critical missing piece in the WAM ecosystem is the lack of joint evaluation metrics that quantify the causal consistency between imagined futures and generated actions. To bridge this gap, future benchmarks must introduce coupled metrics that probe the causal link between visual prediction and physical execution. For instance, \textit{Counterfactual Consistency} could quantify how actions adapt to perturbations in the imagined future, while \textit{Foresight-Conditioned Success} ensures executed trajectories strictly adhere to the generated visual plan rather than relying on spurious correlations. Evaluating the ``intent alignment'' between what the model imagines and what it physically executes is essential for ensuring that actions are genuinely grounded in accurate visual foresight rather than memorized dataset biases.


Beyond individual metrics, the community may benefit from a shared benchmarking framework designed around the core premise of WAMs: one in which world prediction quality and action quality are evaluated jointly and in relation to each other, rather than through separate leaderboards. Such a framework, analogous in ambition to what standard benchmarks have provided for visual recognition or language understanding, could help establish community consensus around what it means for a WAM to be ``working'' in the intended sense—not merely predicting plausible futures, nor merely succeeding at tasks, but generating actions that are demonstrably causally grounded in accurate physical foresight.


\highlightpara{Safety and Reliable Physical Deployment.} WAMs deployed in physical environments introduce safety considerations that go beyond the bounded reactive failures of conventional VLA policies. A model that confidently imagines an incorrect physical future may commit to extended action sequences whose real-world consequences are difficult to interrupt or recover from, potentially causing harm to objects, environments, or people. At the same time, the predictive capacity of WAMs offers a principled opportunity for safety enforcement: world predictions can, in principle, be checked against physical constraints or conservative uncertainty estimates before predicted actions are executed. Realizing this opportunity requires integrating safety verification into the inference pipeline in ways that are computationally tractable and robust to distributional shift, which remains a significant open challenge as WAM capabilities and deployment ambitions continue to expand.


More broadly, the safety challenge highlights an underexplored duality in the WAM paradigm: the same predictive capacity that makes these models potentially more capable than reactive policies also makes their failure modes potentially more consequential. Turning this duality into an advantage means using world predictions not only to guide actions but to verify them before execution. This points toward \textit{prediction-integrated safety}, where uncertainty estimates over imagined futures are treated as first-class inputs to a safety monitor rather than as incidental byproducts. Whether such monitors can be made both computationally tractable and sufficiently robust to distributional shift to be meaningful in open-world deployments remains to be established.



